%! Setting Up --- Learning from Data — Unit 8, Lesson 8.0 — Warm-Up (Answer Key)
\documentclass[10pt]{article}
\usepackage{linalg-article}
\usepackage{linalg-key}   % pulls in -boxes; do NOT also load -boxes
\newcommand{\TallMath}[1]{$\displaystyle #1\rule[-1.4em]{0pt}{3.2em}$}

\begin{document}

\pageheader{Unit 8, Lesson 8.0}{Warm-Up --- Answer Key}
\namedateperiod

\vspace{0.12in}

\begin{notesbox}{Getting Ready: Ideas We'll Use Today}
\small These three ideas --- two from earlier units, one from plain arithmetic --- are the pieces a
neural network is built from. Answer quickly.

\vspace{4pt}
\textbf{1. Matrix times a vector (Unit 1).} Compute $A\mathbf{v}$ for
$A=\begin{bmatrix} 1 & -1 \\ 1 & 1 \end{bmatrix}$ and $\mathbf{v}=\begin{bmatrix} 1 \\ 2 \end{bmatrix}$:
\[
  A\mathbf{v}=\begin{bmatrix} 1 & -1 \\ 1 & 1 \end{bmatrix}\begin{bmatrix} 1 \\ 2 \end{bmatrix}
  = \begin{bmatrix}{\color{keyred}\mathbf{-1}}\\[2pt]{\color{keyred}\mathbf{3}}\end{bmatrix}.
\]

\vspace{4pt}
\textbf{2. How far off? (Unit 4).} Least squares measures error by \emph{squaring} it. A guess of $4$
misses a target of $6$. The error is $4-6={\color{keyred}\mathbf{-2}}$, and its \emph{square} is
${\color{keyred}\mathbf{4}}$. (Squaring makes the error positive and punishes big misses.)

\vspace{4pt}
\textbf{3. The max function (arithmetic).} Compute:
\[
  \max(0,\,5)={\color{keyred}\mathbf{5}},\qquad \max(0,\,-2)={\color{keyred}\mathbf{0}},\qquad \max(0,\,0)={\color{keyred}\mathbf{0}}.
\]
\end{notesbox}

\begin{teachernote}
All three reappear immediately. Item~1 is the \emph{linear step} of every network layer ($A\mathbf{v}$,
Unit~1). Item~2 is the \emph{loss}: least squares squares each error, so the network's ``how wrong am I''
score is a sum of squares (Unit~4). Item~3 is exactly the \textbf{ReLU} $=\max(0,x)$ in disguise --- the
one new function today. Debrief item~3 aloud: $\max(0,x)$ keeps positives and zeroes negatives.
\end{teachernote}

\end{document}
